%%%%%%%%%%%%%%%%%%%%%%%%%%%%%%%%%%%%%%%%%%%%%%%%%%%%%%%%%%%%%%%%%%%%%
%% Manuscript prepared from the ACS "template" model document
%% (J. Wright, CC0). The ACS do not typeset accepted articles with
%% LaTeX, so no journal-specific class is required and this file does
%% not attempt to reproduce the typeset journal layout.
%%%%%%%%%%%%%%%%%%%%%%%%%%%%%%%%%%%%%%%%%%%%%%%%%%%%%%%%%%%%%%%%%%%%%
\documentclass[letterpaper]{article}

%%%%%%%%%%%%%%%%%%%%%%%%%%%%%%%%%%%%%%%%%%%%%%%%%%%%%%%%%%%%%%%%%%%%%
%% Font setup - delete if you are using LuaLaTeX
%%%%%%%%%%%%%%%%%%%%%%%%%%%%%%%%%%%%%%%%%%%%%%%%%%%%%%%%%%%%%%%%%%%%%
\usepackage[T1]{fontenc}

%%%%%%%%%%%%%%%%%%%%%%%%%%%%%%%%%%%%%%%%%%%%%%%%%%%%%%%%%%%%%%%%%%%%%
%% Margins and line spacing
%%%%%%%%%%%%%%%%%%%%%%%%%%%%%%%%%%%%%%%%%%%%%%%%%%%%%%%%%%%%%%%%%%%%%
\usepackage{geometry}
\geometry{margin = 1in}
\usepackage{setspace}

%%%%%%%%%%%%%%%%%%%%%%%%%%%%%%%%%%%%%%%%%%%%%%%%%%%%%%%%%%%%%%%%%%%%%
%% Reference support (biblatex with the ACS chemistry style)
%%%%%%%%%%%%%%%%%%%%%%%%%%%%%%%%%%%%%%%%%%%%%%%%%%%%%%%%%%%%%%%%%%%%%
\usepackage[backend = biber, style = chem-acs]{biblatex}
\addbibresource{acs-main.bib}
\usepackage{xurl}

%%%%%%%%%%%%%%%%%%%%%%%%%%%%%%%%%%%%%%%%%%%%%%%%%%%%%%%%%%%%%%%%%%%%%
%% Graphic inclusion and scheme support
%%%%%%%%%%%%%%%%%%%%%%%%%%%%%%%%%%%%%%%%%%%%%%%%%%%%%%%%%%%%%%%%%%%%%
\usepackage{graphicx}
\usepackage{float}
\newfloat{scheme}{htbp}{los}
\floatname{scheme}{Scheme}

%%%%%%%%%%%%%%%%%%%%%%%%%%%%%%%%%%%%%%%%%%%%%%%%%%%%%%%%%%%%%%%%%%%%%
%% Common support packages
%%%%%%%%%%%%%%%%%%%%%%%%%%%%%%%%%%%%%%%%%%%%%%%%%%%%%%%%%%%%%%%%%%%%%
\usepackage[version = 4]{mhchem}   % formulas via \ce{}
\usepackage{amsmath,amssymb}
\usepackage{booktabs}
\usepackage{siunitx}
\sisetup{group-separator = {,}}

%%%%%%%%%%%%%%%%%%%%%%%%%%%%%%%%%%%%%%%%%%%%%%%%%%%%%%%%%%%%%%%%%%%%%
%% Many journals require that sections are unnumbered
%%%%%%%%%%%%%%%%%%%%%%%%%%%%%%%%%%%%%%%%%%%%%%%%%%%%%%%%%%%%%%%%%%%%%
\setcounter{secnumdepth}{-1}

%%%%%%%%%%%%%%%%%%%%%%%%%%%%%%%%%%%%%%%%%%%%%%%%%%%%%%%%%%%%%%%%%%%%%
%% Author and title data
%%%%%%%%%%%%%%%%%%%%%%%%%%%%%%%%%%%%%%%%%%%%%%%%%%%%%%%%%%%%%%%%%%%%%
\usepackage{authblk}
\author[1,2]{Sunoy Sanyal*}
\author[1,3]{Canberk Demir}
\author[1]{Ali Ozan Aky\"uz}
\author[1]{Sarp D\"onerta\c{s}}
\affil[1]{Department of Chemical and Biological Engineering, Ko\c{c} University, Istanbul, T\"urkiye}
\affil[2]{Department of Industrial Engineering, Ko\c{c} University, Istanbul, T\"urkiye}
\affil[3]{Department of Chemistry, Ko\c{c} University, Istanbul, T\"urkiye}

\title{Closing the Carbon and Nitrogen Loop through Urea Production
  from Nitrate-Rich Wastewater}
% Use the \date command for email address(es) of corresponding authors
\date{*Email: sunoysanyal@gmail.com}

\begin{document}

\maketitle

\begin{abstract}
The United Kingdom faces a coupled environmental challenge in managing
nitrate-rich wastewater and reducing carbon emissions. This study
proposes and evaluates a green urea plant that converts nitrate-rich
brine from Keldgate Water Treatment Works into urea using green
hydrogen and captured biogenic \ce{CO2}. The aim is to close the
nitrogen and carbon loops within a single process and thereby address
both environmental problems simultaneously. The process is modeled
through a combination of mathematically formulated PEM electrolysis,
ASPEN PLUS\textsuperscript{\textregistered} simulation of
Pd--Cu-catalyzed nitrate hydrogenation, ammonia recovery, intermittent
ammonia production, and Bosch--Meiser urea synthesis within a
surrogate-assisted mixed-integer planning framework that accounts for
hourly wind variability, grid-use constraints, and storage design. In
the supporting models, the PEM electrolyzer is assigned a Faradaic
efficiency of $0.98$ and a specific energy demand of
$50~\mathrm{kWh\,kg^{-1}_{H_2}}$. The ammonia-recovery step uses tuned
nitrate-reduction kinetics that reproduce ammonia selectivities of
approximately $85$--$92\%$, together with a membrane-based separation
energy of $7.1~\mathrm{kWh\,kg^{-1}_{N}}$. The unit-level ReLU
surrogates reproduce the simulated plant responses with coefficients
of determination between $0.83$ and $0.999$, with the urea-unit
electricity demand captured at $R^2 = 0.999$. For the small,
DEF-oriented case-study scale of approximately
$125~\mathrm{kg\,d^{-1}}$ urea, all four grid-integration
policies---unrestricted grid, $10\%$-grid, $5\%$-grid, and
wind-only---converge to essentially the same optimal design and
cost: a wind capacity of $1.8$--$1.9$~MW, a $490$~kW electrolyzer, a
$3.2$--$3.5$~MWh battery, and an annualized cost of
\pounds$14.56$--$14.57~\mathrm{million\,y^{-1}}$. Dedicated wind
undercuts grid electricity so decisively that the optimizer purchases
no grid power even when imports are unrestricted, making full
renewable autonomy effectively free at this scale. Intermittency is
absorbed through moderate wind overbuilding with $14$--$16\%$
curtailment, battery and ammonia-inventory buffering, and
intermittent ammonia-section operation at $81\%$ availability, while
fixed plant capital accounts for more than $98\%$ of the total cost.
The resulting break-even urea price of approximately
\pounds$404~\mathrm{kg^{-1}}$, against a market price of
\pounds$1.83~\mathrm{kg^{-1}}$, shows that economic viability is
limited by plant scale rather than by renewable intermittency. Within
the gate-to-gate boundary, with the dedicated wind supply
carbon-neutral at the point of use, the process is a net consumer of
\ce{CO2} at $-0.73~\mathrm{kg_{CO_2}\,kg^{-1}}$ in every case;
including life-cycle electricity intensities as a sensitivity yields
$+1.6~\mathrm{kg_{CO_2}\,kg^{-1}}$, comparable to conventional
production, which identifies the specific electricity demand as the
key maturation lever. These results indicate that nitrate-to-urea
valorization is technically feasible and compatible with fully
islanded wind operation, but that its economics call for feedstock
aggregation toward larger scales or remuneration of the
nitrate-removal service itself, rather than support directed at the
renewable-energy supply.
\end{abstract}

\section*{Keywords}

green urea; nitrate valorization; wastewater; surrogate modeling;
mixed-integer programming; renewable intermittency; carbon utilization

\section*{Abbreviations}

BEP, break-even price; CAPEX, capital expenditure; CRF, capital
recovery factor; DEF, diesel exhaust fluid; FE, Faradaic efficiency;
IX, ion exchange; LCA, life-cycle assessment; MILP, mixed-integer
linear programming; NPW, net present worth; OPEX, operating
expenditure; PEM, proton-exchange membrane; ReLU, rectified linear
unit; SCR, selective catalytic reduction; WTW, water treatment works.

%%%%%%%%%%%%%%%%%%%%%%%%%%%%%%%%%%%%%%%%%%%%%%%%%%%%%%%%%%%%%%%%%%%%%
%% Main text
%%%%%%%%%%%%%%%%%%%%%%%%%%%%%%%%%%%%%%%%%%%%%%%%%%%%%%%%%%%%%%%%%%%%%
\section{Introduction}

The transition toward sustainable chemical manufacturing has
intensified research into alternative nitrogen management strategies
that reduce both environmental impact and reliance on fossil
resources. Conventional urea production is intrinsically linked to
ammonia synthesis via steam methane reforming, a process that is both
energy-intensive and a major contributor to global \ce{CO2}
emissions. A review by Ghavam et al.\ reports a carbon footprint of
2.16~kg\,\ce{CO2}-eq per kg \ce{NH3} for the Haber--Bosch process and
an energy demand of 30~GJ per tonne \ce{NH3}, both of which place a
substantial burden on the environment and the energy system
\cite{Ghavam2021SustainableAmmonia}. In response, emerging process
concepts focus on decarbonizing hydrogen production through water
electrolysis and integrating carbon capture, thereby enabling
lower-emission synthesis routes for nitrogen-based chemicals.

Nitrate contamination in water bodies, primarily driven by
agricultural activities, has become a significant environmental
concern. Globally, shallow groundwater in agricultural areas often
exceeds $10\,\mathrm{mg\,L^{-1}}\;\mathrm{NO_3^{-}{-}N}$
\cite{Nolan2006Nitrate}, creating a particular challenge for the UK,
where groundwater is an important source of drinking water. Water
treatment facilities increasingly employ ion-exchange processes to
remove nitrate from drinking water, generating concentrated
nitrate-rich brine streams as a secondary waste. These streams,
exemplified by UK installations such as Keldgate Water Treatment
Works, contain substantial nitrate concentrations and represent a
potentially valuable nitrogen resource rather than solely a disposal
challenge. The conversion of nitrate-rich wastewater into useful
chemical products has emerged as a promising strategy that
simultaneously addresses pollution mitigation and resource
recovery. This study focuses on the conversion of such waste streams
into urea through integrated chemical design.

The importance of this work lies in its potential to address multiple
challenges simultaneously. First, it provides a pathway for the
conversion of nitrate waste, transforming an environmentally
problematic byproduct into a valuable chemical commodity. In doing
so, it contributes to closing the anthropogenic nitrogen loop,
whereby reactive nitrogen discharged into the environment is
recovered and reintegrated into the chemical value chain. This
approach aligns closely with circular economy principles, which
emphasize resource recovery, waste minimization, and the
recirculation of materials within industrial systems. Second, by
incorporating hydrogen produced via water electrolysis, the proposed
process contributes to the reduction of greenhouse gas emissions
associated with conventional nitrogen chemical production. Third, the
study explores the design and operation of small-scale, distributed
chemical systems, which represent a departure from traditional
large-scale centralized production and may offer advantages in terms
of resource efficiency and process integration.

The main contributions of this work are threefold. It develops a
process design framework based on real wastewater data, evaluates the
suitability of different technological pathways for nitrate-to-urea
conversion, and assesses the economic implications of operating at
sub-industrial scale. In doing so, the study provides insights into
the opportunities and limitations of integrating wastewater treatment
with sustainable, circular chemical manufacturing systems.

\section{Process Concept and Integration in UK Electrical and
  Chemical Networks}

\subsection{Choice of Key Chemical}

Urea is selected as the target product due to its global significance
as a nitrogen carrier and its central role in fertilizer
production. As one of the highest-volume chemicals in the world
\cite{IFASTAT_UreaAtAGlance}, urea serves as a critical link between
closing the nitrogen loop and agricultural productivity. Moreover,
its synthesis from ammonia and carbon dioxide provides a direct route
for carbon utilization, aligning with broader decarbonization
objectives.

From a process integration perspective, urea production offers a
compelling pathway for coupling nitrate recovery with green hydrogen
utilization. Nitrate-rich waste streams can be converted to ammonia
via catalytic hydrogenation, after which conventional urea synthesis
chemistry can be applied. At the scale of typical water treatment
facilities, the nitrate content of waste brine streams is sufficient
to support continuous, albeit small-scale, urea production on the
order of approximately one tonne per day.

Despite the maturity of urea synthesis at industrial scale, its
implementation at smaller scales introduces techno-economic
challenges. These challenges---particularly those associated with
integrating ammonia and urea synthesis and with the high capital
intensity of the process---motivate further investigation into
process design and optimization.

The following subsections outline the integration of this process
into the United Kingdom's network of chemical production, waste
feedstock, and renewable electricity. Key parameters such as the
waste source, carbon dioxide source, and product destination are
considered first. Afterwards, a renewable-energy source is identified
to finalize the proposed site location.

\subsection{Source of Waste}
% THIS PART WAS WRITTEN BY CANBERK DO NOT DELETE
In this work, the waste stream is chosen as nitrate-rich brine
generated within the urban water-supply network. At Keldgate Water
Treatment Works (WTW) in Yorkshire, UK, drinking water is treated
using ion-exchange (IX) columns. When these IX columns are
regenerated, a nitrate-rich brine is produced. This brine contains
high concentrations of chlorine together with smaller amounts of
phosphate, sulphate, and nitrate ions. Among these species, nitrate
is of particular interest because it can be hydrogenated to form
ammonia. The Keldgate plant has a total treatment capacity of
approximately 90~ML\,day$^{-1}$, with a nitrate-removal sidestream
of up to 33~ML\,day$^{-1}$, achieving treated nitrate
concentrations below 5~mg\,L$^{-1}$. Based on reported
operating data \cite{Corrigan2010} and mass-balance estimates, this
corresponds to a brine flowrate in the range of
180--450~m$^3$\,day$^{-1}$, with an average design basis of
207~m$^3$\,day$^{-1}$. From the data provided by Keldgate
WTW, the nitrate concentration in the brine is estimated to be
5~g\,L$^{-1}$. The design basis for Keldgate WTW corresponds to
a nitrate recovery rate of approximately 1.9~t\,day$^{-1}$, which
defines the scale of the plant discussed in this paper.

\subsection{Destination of Product}
The primary product from this process is intended for use as diesel
exhaust fluid (DEF, AUS32), a 33.3~wt\% aqueous urea solution
employed in selective catalytic reduction (SCR) systems for the
abatement of nitrogen oxide (NO$_x$) emissions from diesel
engines. Within the UK, DEF represents a well-established and growing
market, driven by increasingly stringent emissions regulations in
both on-road and off-road transport sectors \cite{Koebel2000}. At the
production scale considered in this work (sub-tonne-per-day urea
output), direct entry into bulk fertilizer markets is economically
unfavorable because the process does not benefit from large-scale
economies. In contrast, DEF is a higher-value product for which
smaller, decentralized production units may be competitive,
particularly when integrated with local supply chains. The production
of DEF from nitrate-derived urea therefore provides a viable pathway
for valorizing waste streams while supplying a product with clear
environmental functionality.

\subsection{Source of Renewables}
In addition to feedstock availability, another critical factor is the
sourcing of low-carbon electricity, particularly for the
electricity-intensive hydrogen-production step via electrolysis,
which dominates the process energy demand. This requires careful
consideration of an appropriate power source. In line with the energy
requirements of the process, the UK is an attractive location because
it hosts some of the world's largest offshore wind facilities in the
North Sea. The Hornsea~2 offshore wind farm has a total capacity of
1.32~GW, provided by 165 turbines of 8~MW each
\cite{OrstedHornsea22026}. Assuming that a portion of this capacity
can be reserved for the present process, the estimated total capital
cost of the facility (\pounds 1.1412~billion)
\cite{OfgemHornsea22023} can be normalized by the installed capacity
of 1320~MW, yielding a specific capital cost of approximately
\pounds 864{,}000 per MW. Applying this value to the wind-capacity
requirement that arises in the case studies (approximately 1.9~MW)
gives an estimated installation cost of about \pounds 1.6~million.

\subsection{Source of Carbon Dioxide}
A critical feedstock for the process is \ce{CO2}, which is consumed
in urea synthesis via the Bosch--Meiser process. This approach
supports UK decarbonization objectives by simultaneously closing both
the nitrogen and carbon loops within the process. Nitrogen is
recovered from wastewater in the form of nitrate and converted into
ammonia and subsequently urea, thereby avoiding its release to the
environment and reducing reliance on conventional nitrogen fixation
routes. In parallel, the urea-synthesis step incorporates carbon
dioxide as a reactant, enabling partial utilization of captured or
waste \ce{CO2} streams. When combined with low-carbon hydrogen, this
integrated pathway offers the potential for significantly lower
life-cycle carbon emissions than conventional fossil-based urea
production. The proposed process uses biogenic \ce{CO2} derived from
anaerobic digestion at the Sandhill Biogas facility, which has a
capacity of 4{,}000~t\ce{CO2}/year, or about 11~t/day
\cite{DESNZTeessideCCUS2025}. Compared with sourcing \ce{CO2} from
industrial flue gases, this approach recycles carbon already present
in the short-term carbon cycle. It is also advantageous because
\ce{CO2} streams from biogas upgrading are generally available at
higher purity, reducing the need for extensive pretreatment. The
proposed \ce{CO2} demand for the process is 0.78~t/day. The Sandhill
Biogas facility is located approximately 23~km from the proposed
site, which further supports its suitability as a feedstock source,
as its permit schedules explicitly incorporate \ce{CO2} conditioning
and liquefaction operations, including capture, compression, cooling,
and condensation, with the resulting liquid \ce{CO2} removed off-site
by tanker transport, making small offtakes feasible.

\subsection{Choice of Location}
An optimal UK site should satisfy several constraints, including
proximity to nitrate-rich wastewater brine, viable access to offshore
wind infrastructure, location within a region with suitable wind and
climate data, and integration with industrial \ce{CO2} transport
infrastructure. When these constraints are considered, the East
Yorkshire/Humber industrial region emerges as an appropriate
candidate site. The region is attractive because of its proximity to
Keldgate WTW, which provides access to nitrate-rich brine feedstock,
its proximity to Hornsea~2 as a renewable-energy source, and its
viable access to \ce{CO2} from the Sandhill Biogas facility via
tanker transport. In addition, its location within an established
industrial cluster provides strong grid connectivity and access to
supporting infrastructure.

\section{Design and Simulation of the Integrated Plant}

\subsection{Overall Plant Design}
% WRITTEN BY CANBERK
To integrate the Haber--Bosch and Bosch--Meiser process elements, a
combination of simulation, mathematical modeling, and Python-based
integration was used. The Bosch--Meiser process was modeled entirely
in ASPEN with slight modifications based on the patent described by
Pustjens \cite{Pustjens2018}. The design of the green ammonia plant
combined mathematical modeling and simulation. The electrolyzer block
and stripping loop were modeled mathematically, whereas the reaction
loop was designed in ASPEN. The overall green ammonia plant was then
embedded in the optimization step through integration of the ASPEN
simulation, mathematically modeled blocks, and intermittency
units. The ELECNRTL fluid package was used for the green ammonia
plant, whereas the SR-POLAR fluid package was used for the
Bosch--Meiser process. The overall process can be divided into five
main parts: the electrolyzer block, ammonia catalytic hydrogenation,
ammonia separation, the urea synthesis loop, the urea separation
section, and intermittency considerations.

\begin{figure}
\centering
\includegraphics[width=\textwidth]{prism-uploads/prodScheme.png}
\caption{Simplified process flow diagram of the plant design showing
  the main subprocesses.}
\label{fig:prod_scheme}
\end{figure}

\subsection{Electrolyzer Block}
In this work, the electrolyzer block was separated from the main
simulation because its scale and intermittent operation are modeled
independently of the green ammonia plant. The mathematical model of
the electrolyzer block was formulated with explicit consideration of
water use and energy demand. Stoichiometry was defined from the
overall reaction, and separate anodic and cathodic reactions were not
considered. In addition, the Faradaic efficiency (FE) of the cell was
incorporated into the mass calculation to determine the actual water
requirement needed to satisfy the hydrogen demand.

\begin{equation}
\ce{H2O(l) -> H2(g) + 1/2 O2(g)}
\end{equation}

\begin{equation}
m_{\ce{H2},\mathrm{actual}} = \frac{m_{\ce{H2},\mathrm{stoich}}}{\eta_F}
\end{equation}

\begin{equation}
m_{\ce{H2O},\mathrm{actual}}
  = \frac{0.5\,m_{\ce{O2},\mathrm{stoich}}}{\eta_F}
  = \frac{m_{\ce{H2},\mathrm{stoich}}}{\eta_F}
\end{equation}

\noindent For the energy demand, a specific-energy scalar was
introduced to convert mass flow rate into electricity demand.

\begin{equation}
E = \frac{m_{\ce{H2}} \cdot e_{\mathrm{spec}}}{\eta_F}
\end{equation}

\noindent The screening of water electrolyzers and the derivation of
$e_{\mathrm{spec}}$ are provided in the Supporting Information. After
screening, FE was taken as 0.98 based on the work of Rashid et
al.\ \cite{Rashid2015}, and $e_{\mathrm{spec}}$ was calculated as
$50\,\mathrm{kWh}\,\mathrm{kg}_{\ce{H2}}^{-1}$.

\subsection{Ammonia Catalytic Hydrogenation}
The nitrate hydrogenation reactor is modeled using a lumped
pseudo-first-order kinetic approach with the catalyst chosen as the
industrial standard Pd--Cu. Thus, a reaction network was introduced
and tuned for the operation of the Pd--Cu catalyst at 60~$^\circ$C
and 1~bar. In this tuning the work of Zhang et
al.\ \cite{Pd-Cu_experimental} was taken as a basis and \ce{NH3}
selectivity of 90\% and \ce{NO3-} conversion of 100\% was aimed. The
proposed kinetic modeling is as follows:

\begin{align}
\ce{NO3- + H2 &-> NO2- + H2O} \\
\ce{NO2- + 3H2 &-> NH3 + 2H2O} \\
\ce{2NO3- + 5H2 &-> N2 + 6H2O} \\
\ce{2NO2- + 3H2 &-> N2 + 4H2O}
\end{align}

\begin{equation}
\begin{aligned}
r_1 &= k_1 C_{\ce{NO3-}} & \qquad r_2 &= k_2 C_{\ce{NO2-}} \\
r_3 &= k_3 C_{\ce{NO3-}} & \qquad r_4 &= k_4 C_{\ce{NO2-}}
\end{aligned}
\end{equation}

\noindent This formulation enables the inclusion of competing pathways
and side reactions. The reaction constants were then tuned by trial
and error, and the final values are reported in the Supporting
Information. To account for intermittent reactor operation, the
autothermal reactor system proposed by Cheema and Krewer
\cite{Cheema2018} was adopted. In the ASPEN simulations, this
autothermal reactor system was modeled as an RPLUG reactor.

\subsection{Ammonia Separation}
Ammonia produced from nitrate hydrogenation exits the reactor as an
aqueous solution (hydrous ammonia), requiring separation and
purification prior to urea synthesis. Because ammonia is highly
miscible in water, this step presents a separation challenge. In
addition, the relatively small process scale limits the practicality
of rigorous stripping system modeling. Accordingly, an alternative
method was adopted: a novel gas stripping membrane reported by Iddya
et al.\ \cite{Iddya2020}. This method has the advantage of separating
ammonia from saline water without side reactions or extreme heat
duties. Following the work of Iddya et al., the specific energy
demand for ammonia recovery was calculated as:

\begin{equation}
Q_S = \dot m_{\ce{NH3}} \cdot e_{\mathrm{membrane}}
\end{equation}

\noindent where $e_{\mathrm{membrane}}$ is specified as
7.1~kWh\,kg$^{-1}$ of nitrogen \cite{Iddya2020}.

\subsection{Urea Reaction}
% WRITTEN BY CANBERK
Urea synthesis is carried out via the conventional reaction between
ammonia and carbon dioxide, proceeding through ammonium carbamate
formation followed by dehydration to urea.

\begin{equation}
\ce{2NH3 + CO2 <=> NH2COONH4}
\end{equation}

\begin{equation}
\ce{NH2COONH4 <=> NH2CONH2 + H2O}
\end{equation}

\noindent The first reaction is fast and exothermic, whereas the
second is slower and endothermic and is therefore both rate and
equilibrium limiting. To favor carbamate formation, the RPLUG reactor
block operates at 200~$^\circ$C and 200~bar, consistent with
industrial practice. Low-pressure decomposition reactors at 20~bar
and 2~bar are then modeled with RSTOIC blocks to decompose excess
carbamate into carbon dioxide and ammonia. Within the reaction loop,
the decomposed gases are recycled back to the reactor. The reaction
kinetics were modeled as surrogates to match the conversion presented
in the work of Pustjens \cite{Pustjens2018}. Again, a lumped
first-order kinetic reaction network was used and tuned by trial and
error. The proposed reaction network consists of four elementary
reactions, representing the two forward and two backward
reactions. The finalized kinetic values can be found in the
Supporting Information.

\begin{align}
\ce{2NH3 + CO2 &-> NH2COONH4} \\
\ce{NH2COONH4 &-> 2NH3 + CO2} \\
\ce{NH2COONH4 &-> NH2CONH2 + H2O} \\
\ce{NH2CONH2 + H2O &-> NH2COONH4}
\end{align}

\subsection{Urea Separation}
The urea-separation cycle consists of two flash drums operating at
4~bar and 1~bar for upstream separation of excess gases. These gases
are then allowed to react to form carbamate, which is simulated using
an RGIBBS reactor. Downstream, the water--urea solution is separated
using a two-stage distillation setup. First, the mixture enters a
stripper, where trace amounts of ammonia are removed and water is
largely vaporized. The remaining urea--water mixture then enters a
distillation column, which vaporizes water from the mixture and
yields an approximately 33~wt\% urea--water solution that can be
adjusted to the exact DEF concentration. The distillation column
operates at 1~bar with a reflux ratio of 1. The number of stages was
kept at a minimum of five. The column returns a total duty of
approximately 0.1~MW, which was considered acceptable and was not
optimized further.

\subsection{Heat Integration and Electric Refrigeration}
% WRITTEN BY CANBERK
Heat integration is a critical component in reducing operating costs
in the nitrate-to-urea process, particularly because of thermally
intensive unit operations such as stripping columns, coolers,
compressors, and the high-pressure, high-temperature urea
reactor. Given the relatively small scale of the plant
($\sim$125~kg/d urea), full industrial-grade heat-recovery networks
are not economically feasible; instead, a simplified but effective
integration strategy is adopted. The primary source of recoverable
heat is the ammonia stripper, which requires significant thermal duty
to remove water from hydrous ammonia. In the surrogate formulation,
this duty is represented as $Q_S$, and partial recovery is
implemented through a heat-recovery term $Q_{\mathrm{rec}}$,
constrained by exchanger capacity:

\begin{equation}
0 \leq Q_{\mathrm{rec}} \leq Q_S
\end{equation}

\noindent Recovered heat is reused internally. The remaining thermal
demand is supplied via external hot utilities:

\begin{equation}
Q_{\mathrm{HU}} = Q_S - Q_{\mathrm{rec}}
\end{equation}

\noindent This formulation ensures that heat integration directly
reduces operating expenditure while remaining compatible with the
surrogate-based optimization framework. Given the absence of
large-scale steam networks and the fully electrified design
philosophy, cooling is provided through electrically driven
refrigeration systems. This is consistent with green ammonia and
hydrogen systems, where refrigeration is typically implemented
through closed-loop chillers powered by electricity rather than
steam-driven absorption systems. Refrigeration loads are implicitly
included in the unit-level electricity demands ($E_N$, $E_U$) within
the surrogate models and therefore contribute to the overall plant
power balance. A hybrid approach combining partial heat recovery from
heat-intensive units with electrically driven heating and cooling
utilities provides a robust and scalable solution.

\section{Process Optimization under Intermittent Renewable Energy
  Supply}

The proposed system seeks the least-cost combination of unit sizes
and hourly operating decisions for a wind-powered urea plant under
intermittent wind supply and continuous urea demand. As discussed
above, the Hornsea~2 wind farm is selected as the renewable-energy
source. To obtain an optimal sizing and operating schedule, a
time-varying IPPS formulation is used \cite{Zhang2017IPPS}. This
section outlines the methodology for integrating intermittent wind
energy and presents the mathematical optimization model; the
resulting designs are analyzed in the Results and Discussion.

\subsection{Renewable Energy Integration}
The use of renewable energy has clear advantages for the
sustainability and viability of the process; however, it also
introduces challenges because renewable generation is inherently
variable and dependent on environmental conditions. In the present
design, electricity is supplied by the Hornsea~2 offshore wind farm
on the Yorkshire coast. Because fluctuations in wind intensity affect
the electricity supply, supplementary strategies are required for
optimal and reliable operation. Possible approaches to mitigate the
effects of fluctuating wind conditions include:
\begin{itemize}
  \item \textbf{Energy Storage Systems:} Batteries and other storage
    technologies can be used to buffer fluctuations in energy supply
    \cite{Rashid2016}. Their integration is crucial because wind
    generation varies with climatic conditions, and intermittent
    process operation benefits from the flexibility provided by
    storage. In the present formulation, battery use is a decision
    variable, stored power is an operating variable, and battery
    capacity is a design variable. At the scale considered here,
    battery integration is feasible.
  \item \textbf{Flexible Electricity Utilization:} Non-renewable
    sources such as grid electricity can be used to balance the
    intermittency of wind power. The combination of wind electricity
    and grid electricity therefore provides operational flexibility
    and makes the system easier to adapt to varying conditions.
  \item \textbf{Real-Time Monitoring:} The inclusion of numerous
    sensors and advanced metering infrastructure enables continuous
    data collection and supervision \cite{GaoSheng2018}. Such
    monitoring supports flexible control and more efficient use of
    energy, thereby allowing real-time process adaptation to
    intermittent wind supply.
  \item \textbf{Demand-Side Management:} Active adjustments to
    process operation, enabled by real-time monitoring, can
    synchronize plant activity with periods of high renewable-energy
    availability. This can be achieved by dynamically scaling the
    production rate, and therefore the electricity demand, based on
    instantaneous energy availability.
\end{itemize}

\subsection{Optimization Model}
The proposed system is formulated as a Mixed Integer Linear
Programming (MILP) model. The model jointly decides how large to
build system components, how to operate them each hour, and how to
store and dispatch electricity, hydrogen, \ce{CO2}, and anhydrous
ammonia.

The objective is to minimize the annualized capital cost and
operating costs while satisfying the total demand for urea production
under intermittent wind energy:

\begin{equation}
\min Z = CAPEX_{\mathrm{annualized}} + \sum_{t \in \mathcal{T}} \Delta t \, OPEX_t
\end{equation}
\begin{equation}
\sum_{t \in \mathcal{T}}F^U_t \ge D^{ann}\left(\frac{|\mathcal{T}|}{8760}\right)
\end{equation}

\noindent where $|\mathcal{T}|$ denotes the number of modeled
periods, $F^U_t$ is the pure urea flowrate at a given timestep (hour)
and $D^{ann}$ is the annual pure-urea production target.

The ammonia and urea production units are embedded into the
mathematical model using surrogate-based process models:

\begin{equation}
\big(F^{1,op}, F^{2,op}, E^{op}\big) = \mathcal{S}(\mathbf{x})
\end{equation}

\noindent where $\mathcal{S}(\cdot)$ is a trained ReLU neural-network
mapping. ReLU networks are selected because their piecewise-linear
activation structure admits an exact mixed-integer linear
reformulation, allowing the surrogate models to be embedded directly
within the planning MILP rather than evaluated externally
\cite{Grimstad2019ReLU,Anderson2020Strong}. In this work, the trained
networks are reformulated using the \texttt{reluMIP} library, which
automatically generates MILP encodings of trained ReLU networks and
strengthens the resulting formulation through progressive bound
tightening \cite{Lueg2021reluMIP}. The planning model uses one
time-invariant operating point for each process unit. Thus, the
optimizer chooses constant operating-point inputs and outputs for
both units.

\subsubsection{Energy Storage Dynamics}
To absorb renewable intermittency, the planning model includes
battery, hydrogen, carbon dioxide, and ammonia storage. Only
representative constraints are shown here. The hydrogen and ammonia
balances illustrate the role of storage as a temporal buffer between
production and consumption:
\begin{equation}
\begin{aligned}
F^{H2}_t + ch^H_t &= M^{H2}_t + dis^H_t,\\
F^{NH3}_t + ch^{NH3}_t &= F^{strip}_t + dis^{NH3}_t
\end{aligned}
\end{equation}
Thus, material generated in one period can either be consumed
immediately or shifted to future periods through storage.

\subsubsection{State-of-Charge (SOC) Dynamics}
Inter-temporal consistency is enforced through cyclic inventory
balances. For the battery, the state of charge evolves according to
\begin{equation}
I^B_t = I^B_{t-1} + \eta_B^{ch} ch^B_t - \frac{dis^B_t}{\eta_B^{dis}}
\end{equation}
with cyclic closure $I^B_0 = I^B_T$. Analogous state equations are
imposed for hydrogen, carbon dioxide, and ammonia inventories,
ensuring that the model does not create or deplete storage
nonphysically over the horizon.

\subsubsection{Power-to-Energy Sizing Constraint}
Storage design and operation are linked through upper bounds on
inventory and charging/discharging decisions. For the battery, a
compact representation is adopted:
\begin{equation}
\begin{aligned}
I^B_t &\le B^{cap},\\
ch^B_t &\le P_B^{max} z^B_t,\\
dis^B_t &\le P_B^{max}(1-z^B_t)
\end{aligned}
\end{equation}
where the binary variable $z^B_t$ prevents simultaneous charging and
discharging. In this way, the model captures the main trade-off
between installed storage capacity and short-term dispatch
flexibility without overburdening the planning formulation with
unnecessary detail. The model is implemented in Python using the
Pyomo optimization modeling framework and is solved using the Gurobi
optimizer to a relative optimality tolerance of 1\%
\cite{Hart2011,gurobi}; cost differences between cases that fall
within this tolerance are treated as insignificant in the analysis
below. The wind capacity factor profile,
$CF_t$, is obtained from Renewables.ninja for a wind park with
installed capacity $W^{cap}$ \cite{renewables_ninja}. Accordingly,
the wind power that can be supplied at any given time $t$ from the
installed capacity is bounded by the product of the instantaneous
capacity factor and the installed capacity, leading to the following
availability constraint:
\begin{equation}
  P^{wind}_t \le CF_t W^{cap}
\end{equation}
In addition, grid imports are regulated through a case-dependent
grid-share parameter, $\gamma^{grid}$, such that
\begin{equation}
  P^{grid}_t \le \gamma^{grid} P^{load}_t
\end{equation}
where $\gamma^{grid}$ changes across the four cases and takes the
values 1, 0.10, 0.05, and 0 for the unrestricted-grid, 10\%
grid-limit, 5\% grid-limit, and wind-only cases, respectively. The
simulation accounts for site location, historical weather data, and
turbine height to estimate hourly electricity availability over the
year. Our simulated wind park is the Hornsea~2 cluster, as described
above. The planning horizon comprises 888 hourly periods, and all
capital and operating costs are annualized to a full year through the
factor $8760/|\mathcal{T}|$.

\section{Results and Discussion}

\subsection{Surrogate Accuracy and MILP Embedding}

The fidelity of the planning results rests on how faithfully the
ReLU surrogates reproduce the underlying ASPEN flowsheets, so the
surrogate quality is examined first. Table~\ref{tbl:surrogates}
summarizes the test metrics of the six unit-level outputs used in the
planning model. Figure~\ref{fig:urea-relu-surface} shows the
surrogate data points together with the fitted ReLU neural-network
response surface for the electricity required to produce urea over
the admissible \ce{NH3} and \ce{CO2} operating region.

\begin{table}
  \caption{Test accuracy of the ReLU neural-network surrogates
    embedded in the planning MILP. MAE is reported in the units of
    the respective output.}
  \label{tbl:surrogates}
  \centering
  \begin{tabular}{lllcc}
    \toprule
    Unit & Output & Unit basis & $R^2$ & MAE \\
    \midrule
    Ammonia & \ce{NH3} product   & kg\,h$^{-1}$ & 0.960 & 0.27 \\
    Ammonia & water product      & kg\,h$^{-1}$ & 0.910 & 60.9 \\
    Ammonia & electricity demand & kW           & 0.831 & 185  \\
    Urea    & pure urea product  & kg\,h$^{-1}$ & 0.973 & 0.20 \\
    Urea    & product urea mass fraction & --   & 0.931 & 0.023 \\
    Urea    & electricity demand & kW           & 0.999 & 0.12 \\
    \bottomrule
  \end{tabular}
\end{table}

\begin{figure}
  \centering
  \includegraphics[width=0.75\textwidth]{figures/surface_vs_datapoints.pdf}
  \caption{Comparison of the sampled surrogate data and the fitted
    ReLU neural-network response surface for the electricity
    requirement of the urea unit as a function of \ce{NH3} and
    \ce{CO2} feed rates. The markers represent sampled process data,
    while the surface represents the corresponding surrogate
    prediction. The figure illustrates the strongly nonlinear
    dependence of electricity demand on the feed conditions over the
    admissible operating region. The fitted surrogate achieves a
    coefficient of determination of $R^2 = 0.9985$.}
  \label{fig:urea-relu-surface}
\end{figure}

Two features of Table~\ref{tbl:surrogates} shape the interpretation
of the planning results. First, the urea-unit electricity demand is
captured almost exactly ($R^2 = 0.999$, MAE $=0.12$~kW): over the
admissible feed window the compression, decomposition, and separation
duties respond smoothly to the \ce{NH3} and \ce{CO2} feed rates, so a
compact piecewise-linear network suffices. Second, the ammonia-unit
electricity demand is the least accurate output ($R^2 = 0.831$, MAE
$=185$~kW). This output aggregates the hydrogenation preheater,
cooler, reactor, and compressor duties, whose interaction with the
brine and hydrogen feeds is markedly less regular across the sampled
grid. Because the ammonia unit is also the largest electricity
consumer in the plant (about two thirds of the combined
ammonia--urea process load, cf.\ the energy analysis below), this is
the surrogate error most likely to propagate into the sizing
decisions. The fixed-operating-point formulation mitigates this risk:
the optimizer selects a single operating point per unit, so surrogate
error enters the planning model as a bias at one point of the
response surface rather than as noise across a trajectory, and the
selected operating points can be verified a posteriori against the
rigorous flowsheet. As an additional robustness check, sparse
piecewise-linear hyperplane-tree surrogates with ridge and
elastic-net leaf regressions were trained on the same data and
achieved $R^2$ values between $0.95$ and $0.9997$; a full comparison
is provided in the Supporting Information.

\subsection{Optimal Design and Operating Policy across
  Grid-Integration Cases}

Four operating-policy cases are considered in order to evaluate how
progressively tighter restrictions on grid-electricity utilization
affect plant sizing, storage requirements, and operating
flexibility: an unrestricted-grid benchmark
($\gamma^{grid}=1$), moderately and strictly constrained transition
scenarios in which cumulative grid imports may cover at most 10\% and
5\% of the total plant electricity demand, and a fully islanded
wind-only case in which grid imports are prohibited and all
variability must be absorbed by wind generation, process flexibility,
and storage. In all cases the high-pressure urea section is
constrained to operate continuously, reflecting the intolerance of
the synthesis loop to on/off cycling, so intermittency must be
absorbed upstream of the ammonia buffer. Table~\ref{tbl:cases}
compares the optimal designs, and
Figure~\ref{fig:electricity-balance-cases} shows representative
zoomed-in electricity-balance intervals that illustrate the resulting
short-term operating behavior.

\begin{table}
  \caption{Optimal design capacities and operating characteristics
    for the four grid-integration cases. Availability denotes the
    fraction of hours in which the ammonia section operates; the grid
    share is the realized fraction of total electricity demand
    supplied by the grid, and curtailment is the fraction of
    available wind generation that is spilled over the modeled
    horizon. All four solutions are optimal within a 0.2\% relative
    gap.}
  \label{tbl:cases}
  \centering
  \begin{tabular}{lcccc}
    \toprule
    & Unrestricted & 10\% grid & 5\% grid & Wind only \\
    \midrule
    Wind capacity (kW)            & 1848  & 1813  & 1842  & 1855 \\
    Electrolyzer capacity (kW)    & 490   & 490   & 490   & 491 \\
    Battery capacity (kWh)        & 3212  & 3547  & 3280  & 3150 \\
    \ce{H2} storage (kg)          & 22    & 21    & 18    & 15 \\
    \ce{NH3} storage (kg)         & 477   & 482   & 487   & 483 \\
    \ce{H2} feed at operating point (kg\,h$^{-1}$) & 9.61 & 9.61 & 9.61 & 9.62 \\
    Ammonia-unit availability (\%) & 81.4 & 81.4  & 81.4  & 81.4 \\
    Realized grid share (\%)      & 0     & 0     & 0     & 0 \\
    Wind curtailment (\%)         & 16.1  & 14.0  & 15.9  & 15.8 \\
    Annualized cost (M\pounds\,y$^{-1}$) & 14.565 & 14.565 & 14.565 & 14.565 \\
    \bottomrule
  \end{tabular}
\end{table}

\begin{figure}
  \centering
  \includegraphics[width=\textwidth]{figures/electricity_balance_cases.pdf}
  \caption{Zoomed-in electricity-balance profiles for the four
    grid-integration cases: (a) unrestricted grid, (b) 10\% grid
    limit, (c) 5\% grid limit, and (d) wind only. Bars above zero are
    available electricity supply, bars below zero are electricity
    use, and the dashed line is the battery inventory on the right
    axis; the excess of the supply bars over the use bars corresponds
    to curtailed wind generation.}
  \label{fig:electricity-balance-cases}
\end{figure}

The single most important observation in Table~\ref{tbl:cases} is
that the grid-integration policy has no effect on the optimal
design: all four cases converge to the same solution---a
1.81--1.86~MW wind farm, a 490~kW electrolyzer, a 3.2--3.5~MWh
battery, an ammonia inventory of about 480~kg, and an annualized
cost of \pounds 14.565~million\,y$^{-1}$ that agrees across cases to
within \pounds 500\,y$^{-1}$ (relative gaps 0.16--0.20\%). The
reason is visible in the realized grid shares: even when grid
imports are unrestricted, the optimizer purchases no grid
electricity at all. At the Hornsea~2 specific capital cost of
\pounds 864.5 per kW and the annualization factor of the model, wind
electricity costs roughly \pounds 8~MWh$^{-1}$ at the site mean
capacity factor of 0.594---more than an order of magnitude below the
assumed grid tariff of \pounds 180~MWh$^{-1}$---so the grid
constraint that distinguishes the cases is never active. Renewable
autonomy is therefore not a cost constraint at this scale; it is the
cost-optimal operating mode, and prohibiting grid use entirely is
free.

The common design absorbs wind variability through three
complementary mechanisms. First, generation is moderately overbuilt:
the 1.85~MW farm serves a mean electrical load of about 1~MW, and
14--16\% of the available wind energy is spilled, because at
\pounds 43 per kW-year an occasionally idle turbine is cheaper than
additional storage. Second, the ammonia section operates
intermittently at 81.4\% availability with a deliberately oversized
operating point (design hydrogen feed 9.61~kg\,h$^{-1}$, about 23\%
above the 7.8~kg\,h$^{-1}$ continuous-equivalent requirement),
overproducing ammonia during windy hours and drawing down the
roughly 480~kg ammonia inventory during lulls while the urea section
runs continuously. Third, the battery smooths sub-daily
fluctuations, and a small hydrogen store (15--22~kg) covers the
residual electrolyzer-side variability. The small remaining spread
of the storage capacities across the four columns of
Table~\ref{tbl:cases} carries no design information---it is residual
degeneracy within the 0.2\% optimality gap, as near-optimal designs
can trade battery, inventory, and spill against one another at
almost no cost. Internal heat recovery covers 90--99\% of the
stripper duty $Q_S$, with the balance supplied electrically during
hours of surplus wind.

\subsection{Storage Sizing Hierarchy}

Figure~\ref{fig:storage-comparison} compares the storage capacities
across the four cases. All designs combine a battery of
3.2--3.5~MWh, an ammonia inventory of 477--487~kg, and a
comparatively small hydrogen store of 15--22~kg, and the merit order
behind this mix is consistent. The battery is the preferred
short-term buffer because it is the cheapest and most efficient:
there is little loss of energy and no additional conversion penalty
on charging and discharging. Ammonia inventory is the preferred
\emph{chemical} buffer: whereas hydrogen must be stored under
cryogenic or high-pressure conditions (\pounds 645.89 per kg of
capacity in the cost model), ammonia storage takes place at
near-ambient conditions (\pounds 1.89 per kg) and holds a privileged
position in the flowsheet---discharging it feeds the urea section
directly, without re-running the electrolysis and hydrogenation
steps---which is precisely what sustains the continuous urea
operation required in every case. Hydrogen storage, the only carrier
able to time-shift the electrolyzer load itself, is deployed only
sparingly (15--22~kg).

\begin{figure}
  \centering
  \includegraphics[width=0.8\textwidth]{figures/storage_capacity_cases.pdf}
  \caption{Comparison of storage capacities across the investigated
    system-design cases: hydrogen and ammonia storage (left axis, kg)
    and battery capacity (right axis, kWh).}
  \label{fig:storage-comparison}
\end{figure}

A second, less conventional element of the storage strategy is
generation overbuild. Because a kilowatt of wind capacity costs only
about \pounds 43\,y$^{-1}$ after annualization, installing surplus
turbine capacity and spilling its output during windy hours is
frequently cheaper than cycling energy through any storage medium.
The optimizer exploits this in every case---curtailment of 14--16\%
(Table~\ref{tbl:cases})---so storage is sized only for the residual
variability that overbuild cannot remove, namely sustained low-wind
periods. The resulting design rule is that overbuilt generation
absorbs the bulk of the hourly variability, the battery handles
short-term fluctuations, ammonia inventory smooths multi-hour
production lulls in front of the always-on urea section, and
hydrogen storage covers the deepest electrolyzer-side deficits. None
of the storage capacities approaches its engineering bound (the
battery, for example, uses at most 3.5 of an allowed 20~MWh),
confirming that the designs are cost-shaped rather than
constraint-shaped.

\subsection{Cost Structure and the Cost of Renewable Autonomy}

Figure~\ref{fig:annualized-cost} decomposes the annualized cost of
the four designs by component. The decomposition shows that the cost
of the plant is, to first order, the cost of the chemical plant
itself: the fixed capital of the ammonia section
(\pounds 191.7~million) and the urea section (\pounds 95.8~million)
together account for more than 98\% of the total capital base of
approximately \pounds 291.5~million, and hence of the annualized cost
of \pounds 14.57~million\,y$^{-1}$ at the capital recovery factor of
0.05. The entire renewable-energy system---wind farm
(\pounds 1.6~million), electrolyzer (\pounds 2.0~million),
battery (\pounds 0.33--0.37~million), and chemical
storage (under \pounds 0.02~million)---adds up to roughly 1.4\% of
the capital cost, and the net operating cost is negligible because
the oxygen by-product credit approximately offsets the small
\ce{CO2} feed expenditure.

\begin{figure}
  \centering
  \includegraphics[width=0.9\textwidth]{figures/annualized_cost_allocation.pdf}
  \caption{Allocation of the annualized cost by capital and operating
    component for the four grid-integration cases.}
  \label{fig:annualized-cost}
\end{figure}

Two consequences follow. First, the price of renewable autonomy is
zero: prohibiting grid imports does not change the optimal design at
all, because the wind-plus-storage system that full autonomy
requires is cheap relative to the plant it powers and is fully built
even in the unrestricted-grid optimum. For policy
purposes, the economically meaningful barrier to deploying this
process is not the renewable-energy supply---which the optimization
selects voluntarily on cost grounds---but the capital intensity of
small-scale chemical synthesis. Second, the cost analysis identifies
where effort should be directed: reducing the fixed capital of the
ammonia and urea sections (through simplification, modularization,
or scale-up) has two orders of magnitude more leverage than any
optimization of the energy system, whose design freedom is worth at
most a few percent of the total cost.

\subsection{Energy-Demand Structure}

The distribution of the electrical load between the process sections
is stable across cases: the electrolyzer accounts for approximately
40\% of the total electricity demand, the ammonia section
(hydrogenation, recovery, and compression duties) for a further
35\%, and the urea section for about 16\%, with the remainder
consumed by the electric hot utility, storage auxiliaries, and
battery charging. The urea-section load is constant at 161~kW owing
to its continuous operation. In aggregate the plant consumes
approximately 190~kWh of electricity per kilogram of urea
produced---one to two orders of magnitude above the specific energy
demand of integrated fossil-based ammonia--urea complexes---which is
the single number that most clearly characterizes the efficiency
challenge of the present flowsheet. Two deliberately conservative
conventions inflate this figure. First, the unit-level demands are
electricity equivalents formed from the sum of the absolute
equipment duties, so cooling and condenser loads are charged at a
coefficient of performance of unity: at the ammonia operating point,
the exothermic reactor duty and the effluent cooler account for
roughly 380 of the 460~kW-equivalent, although the temperature level
(60~$^\circ$C) would in practice allow cooling water rather than
electrically driven refrigeration. Second, the hydrogenation is fed
roughly twice the stoichiometric hydrogen requirement (a design feed
of 9.6~kg\,h$^{-1}$ against a stoichiometric demand of about
4.4~kg\,h$^{-1}$ for the produced ammonia) and the unreacted excess
is not recycled, so electrolysis alone contributes about
77~kWh\,kg$^{-1}$ of urea. Figure~\ref{fig:energy-breakdown} resolves the
wind-only case down to individual equipment duties. Within the
ammonia section, the separation and thermal-conditioning duties
(stripper utility, preheater, cooler, and compressor) dominate, while
within the urea section the reboiler and condenser duties of the
separation columns outweigh the synthesis-loop compression. Together
with the electrolyzer share, this breakdown identifies two
high-leverage efficiency targets: the excess-hydrogen operation of
the nitrate hydrogenation (the design hydrogen feed of about
9.5~kg\,h$^{-1}$ far exceeds the stoichiometric requirement of the
recovered nitrate) and the unmatched hydrogenation preheat/cool
pairing in the ammonia thermal train.

\begin{figure}
  \centering
  \includegraphics[width=0.8\textwidth]{figures/wind_only_detailed_energy_duty_breakdown.pdf}
  \caption{Detailed breakdown of the electricity and duty demand for
    the wind-only case, resolved by process section (inner ring) and
    individual equipment duty (outer ring).}
  \label{fig:energy-breakdown}
\end{figure}

\subsection{Economic Evaluation}

The economic evaluation of the system is performed by calculating the
break-even price (BEP) of the urea product using the net present
worth (NPW) framework and comparing the calculated price with current
market prices. This evaluation establishes a basis for assessing
process feasibility. The NPW calculation can be written as
\cite{Peters2025PlantDesign}:

\begin{equation}
\begin{aligned}
\mathrm{NPW} ={}& \sum_{j=1}^{N} PWF_{cf,j}
\left[(s_j - c_{oj} - d_j)(1-\Phi) + rec_j + d_j\right] \\
& - \sum_{j=1}^{N} PWF_{v,j} T_j
\end{aligned}
\end{equation}

\noindent where:
\begin{itemize}
 \item $PWF_{cf,j}$ is the present worth factor for cash flow in year $j$
  \item $s_j$ is the sales revenue in year $j$
  \item $c_{oj}$ is the operating cost in year $j$
  \item $d_j$ is the depreciation in year $j$, assumed to be linear:
    $d_j = (\mathrm{FCI} - C_S) / N$
  \item $\Phi$ is the tax rate
  \item $rec_j$ is the money recovered from the capital investment
  \item $PWF_{v,j}$ is the present worth factor for investment made
    in year $j$, taken to be 1 since the only investments are made in
    the $0^{\mathrm{th}}$ year
  \item $T_j$ is the total investment made in year $j$
    \cite{Peters2025PlantDesign}
\end{itemize}

After substitution of the relevant terms, the NPW formulation can be
simplified as:

\begin{equation}
\begin{aligned}
\mathrm{NPW} ={}& - CAPEX \\
& + \sum_{j=1}^{N} (s_j - OPEX)(1-\Phi)/(1+i)^j \\
& + \sum_{j=1}^{N} d_j/(1+i)^j \\
& + (C_S + C_W)/(1+i)^N
\end{aligned}
\end{equation}

where:

\begin{itemize}
 \item $i$ is the interest rate
  \item $C_S$ is the salvage value of the Fixed Capital Investment
    (FCI), assumed to be $C_S = 0.15 \cdot FCI$; where
    $FCI = 0.80 \cdot CAPEX$
  \item $C_W$ is the working capital, assumed to be
    $C_W = 0.20 \cdot CAPEX$
\end{itemize}

The exact formulations for CAPEX and OPEX can be found in the
Supporting Information. The CAPEX formulation includes all fixed and
working-capital investments associated with the production process,
as well as costs related to wind-farm leasing. The OPEX formulation
includes operational expenses such as grid electricity, raw
materials, maintenance, and general administrative expenses from the
first year of production onward. The evaluation uses an interest
rate of $i = 7\%$, a project lifetime of $N = 10$ years, and a tax
rate of $\Phi = 19\%$ (the UK corporation-tax rate), and values the
revenue streams at the current market prices of urea and \ce{CO2}
certificates. The BEP of urea can then be found by setting NPW equal
to zero and solving the resulting equation for sales revenue under
the four cases.

Because the four designs share the same capital and operating bases,
their economic indicators coincide: the NPW at market prices is
$-104.2$~million \textsterling{} for all four
cases. The negative NPW values indicate that government incentives
would be required to make the process economically feasible. In this
context, the UK government's \pounds 1.3~billion Net Zero Innovation
Portfolio fund for low-carbon systems may represent an important
opportunity for project support \cite{GovUKNZIP2025}: the full NPW
gap of a demonstration plant of this kind amounts to about 8\% of
that fund. It must be emphasized, however, that such support would
subsidize the scale penalty of the chemical plant, not the renewable
energy supply, which the cost-optimal design adopts on its own.

To assess the pricing conditions under which the process would break
even, we calculated the BEP for the two major revenue streams: the
sale of urea and credits from \ce{CO2} certificates. The BEPs of urea
for the different cases are shown in Figure~\ref{fig:bep}, and the
BEPs of \ce{CO2} certificates are shown in
Figure~\ref{fig:bep_cert}. For comparison, the current market price
of urea is 1.83~\textsterling/kg urea, and the current market price
of \ce{CO2} certificates is 41.84~\textsterling/t\ce{CO2}
\cite{GovUKETSCarbonPrice2025}. The computed break-even urea price is
approximately 404~\textsterling/kg in every case---more than two
orders of magnitude above the market price---and the break-even
\ce{CO2}-certificate price (with urea sold at market price and the
credit applied to the roughly 438~t\,y$^{-1}$ of biogenic \ce{CO2}
utilized) is approximately 42{,}000~\textsterling/t, about a
thousandfold the current certificate price. The near-identity of the
break-even prices across the four cases is a direct structural
consequence of the cost decomposition: the ammonia- and urea-plant
capital costs enter through capacity correlations whose fixed terms
are invariant across cases and exceed every case-dependent item by
two orders of magnitude, so the grid policy cannot move the
break-even price by more than a fraction of a percent. It should be
noted that these correlations are extrapolated far below the
capacity range over which they were fitted, so the absolute level of
the break-even price carries substantial uncertainty; the
qualitative conclusion, however---that plant capital rather than the
energy system sets the economics---is robust to this uncertainty.
The gap has a single structural cause: scale. At 125~kg\,d$^{-1}$ the annual product mass
over which the fixed capital is amortized is roughly four orders of
magnitude smaller than in a world-scale urea train, while the
synthesis-plant capital does not scale down proportionally; the
energy system, by contrast, contributes only about 2\% of the capital
base and is economically irrelevant to the gap. Consequently, the
product revenue itself is nearly irrelevant to project economics at
this scale---closing the gap requires either aggregation of feedstock
from multiple treatment works into a larger plant or remuneration of
the nitrate-removal service itself (avoided brine disposal), rather
than support directed at electricity prices or carbon credits.

\begin{figure}
  \centering
  \includegraphics[width=0.7\textwidth]{figures/bep.pdf}
  \caption{Break-even price of urea in each case compared to the
    current market price.}
  \label{fig:bep}
\end{figure}

\begin{figure}
  \centering
  \includegraphics[width=0.7\textwidth]{figures/bep_cert.pdf}
  \caption{Break-even price of \ce{CO2} certificates in each case
    compared to the current market price.}
  \label{fig:bep_cert}
\end{figure}

\subsection{Ecological Evaluation}

A gate-to-gate life-cycle assessment (LCA) is performed to evaluate
the ecological impact of the process. Because urea can be used in a
wide range of downstream products, product-use effects are excluded
from the present LCA. It should nevertheless be acknowledged that a
negative carbon-emission value in this analysis does not necessarily
imply long-term removal of carbon from the atmosphere. Even so,
comparing the different cases remains valuable because all cases are
evaluated using the same metric. In the LCA, avoided greenhouse gas
(GHG) emissions associated with preventing \ce{CO2} release from the
process outlet streams are taken into account. In addition, GHG
emissions associated with electricity use are included. The
functional unit is 1~kg of urea product. The entire \ce{CO2} stream
entering the process is treated as \ce{CO2} removed from the
atmosphere. Besides the carbon that ends up in the urea product,
carbon has two other possible destinations: it may remain unreacted
and be vented to the atmosphere, or it may be combusted and released
as \ce{CO2}. Accordingly, all \ce{CO2} not incorporated into the urea
product is assumed to be released back to the atmosphere.
Electricity-related GHG emissions are divided into two components
because of the hybrid availability of grid electricity and dedicated
offshore wind electricity, with the respective intensities applied to
the electricity actually consumed by the plant (curtailed wind
generation is not charged to the product). Grid electricity carries
the UK operational intensity of 177~kg\ce{CO2}e/MWh
\cite{GovUKGHGConversionFactors2025}. Electricity from the dedicated
wind capacity is carbon-neutral at the point of use, since the
embodied emissions of the turbines lie upstream of the gate-to-gate
boundary; as a sensitivity, the assessment is repeated with the
life-cycle intensity of offshore wind electricity,
13~kg\ce{CO2}e/MWh according to the United Nations Economic
Commission for Europe \cite{UNECECarbonNeutralityLCA2022}. Under
these conventions the material carbon balance contributes a fixed
credit of $44/60 \approx 0.73$~kg\ce{CO2} per kg of urea (the
\ce{CO2} incorporated in the product), and the electricity term
contributes in proportion to the specific electricity demand of
about 190~kWh\,kg$^{-1}$. Transportation-related GHG emissions are
neglected in this assessment. The LCA is performed for all four
grid-integration cases, and the results are shown in
Figure~\ref{fig:emmision}.

\begin{figure}
  \centering
  \includegraphics[width=0.85\textwidth]{figures/emissions.pdf}
  \caption{\ce{CO2} emissions per kilogram of product for the four
    grid-integration cases, in the gate-to-gate accounting with
    dedicated wind carbon-neutral at the point of use and in the
    sensitivity including life-cycle electricity intensities,
    compared with literature values for the conventional production
    route (steam-methane-reforming-based Haber--Bosch ammonia and
    Bosch--Meiser urea).}
  \label{fig:emmision}
\end{figure}

Because every optimal design is fully wind-powered, the gate-to-gate
result is identical across the four cases: the process is a net
consumer of \ce{CO2} at $-0.73$~kg\ce{CO2}\,kg$^{-1}$, the carbon
fixed in the product. Valorizing the nitrate brine therefore yields
urea whose gate-to-gate footprint is negative wherever the dedicated
renewable supply can be treated as carbon-neutral, in contrast to
the conventional fossil-based route at
1.83--2.2~kg\ce{CO2}\,kg$^{-1}$. The sensitivity analysis, however,
shows how strongly this conclusion depends on the electricity
accounting: charging wind electricity with its life-cycle intensity
of 13~kg\ce{CO2}e/MWh raises the intensity to
$+1.56$ to $+1.58$~kg\ce{CO2}\,kg$^{-1}$---comparable to the
conventional benchmarks---because the specific electricity demand of
about 190~kWh\,kg$^{-1}$ then contributes roughly
2.5~kg\ce{CO2}\,kg$^{-1}$, overwhelming the product credit. Had the
plant been grid-powered at the same demand, the intensity would
exceed 30~kg\ce{CO2}\,kg$^{-1}$, so renewable supply remains a
necessary condition in either accounting. The life-cycle sensitivity
also defines a clear maturation target: under that stricter
accounting the process becomes a net remover of \ce{CO2} below a
specific electricity demand of
$0.73/0.013 \approx 56$~kWh\,kg$^{-1}$. The current
190~kWh\,kg$^{-1}$ comprises about 77~kWh\,kg$^{-1}$ of electrolysis
at twice the stoichiometric hydrogen feed and a large cooling
contribution charged at unit coefficient of performance, so the
threshold is plausibly reachable: recycling the excess hydrogen
would cut the electrolysis term toward $\sim$20~kWh\,kg$^{-1}$, and
assigning the near-ambient cooling duties to cooling water rather
than electricity-equivalent refrigeration would remove most of the
remaining gap. Robust carbon benefit therefore hinges on flowsheet
maturation---hydrogen recycle and a conventional utility
system---rather than on the energy supply, which is already fully
renewable in every optimized design.

\section{Conclusions}

This study presents a novel, integrated process for the production of
urea from nitrate-rich wastewater, green hydrogen, and biogenic
carbon dioxide utilization. The proposed process valorizes nitrate by
hydrogenation with green hydrogen, then uses waste carbon dioxide to
produce green urea using a low-pressure Bosch--Meiser process. Thus,
a viable pathway is proposed toward simultaneously addressing nitrate
pollution and carbon emissions, thereby contributing to the closure
of both the nitrogen and carbon loops. The results indicate that,
from a process integration perspective, the concept is technically
feasible and can be effectively synthesized through a hybrid
framework combining first-principles simulation, surrogate modeling,
and mixed-integer optimization.

The incorporation of intermittent renewable energy, particularly
offshore wind, turns out to be the \emph{least} restrictive element
of the design. All four grid-integration policies converge to the
same optimal design and the same annualized cost of
\pounds 14.57~million\,y$^{-1}$, because dedicated wind capacity is
so much cheaper than grid electricity that the optimizer purchases no
grid power even when imports are unrestricted; full grid independence
is free. The case analysis shows a consistent
intermittency-management strategy: moderate wind overbuilding with
14--16\% curtailment absorbs most of the hourly variability, a
battery of about 3.5~MWh handles short-term fluctuations, an ammonia
inventory of about 480~kg smooths production lulls in front of the
continuously operating high-pressure urea section, and a small
hydrogen store covers the deepest electrolyzer-side deficits.

From an economic standpoint, all evaluated scenarios yield negative
NPW ($-$\pounds 104~million at market prices),
indicating that the process is not competitive under current market
conditions without external financial support. The analysis
localizes the gap precisely: more than 98\% of the capital base is
the fixed cost of the small-scale ammonia and urea sections, so the
break-even urea price of about \pounds 404~kg$^{-1}$ reflects a scale
penalty, not an energy penalty. Policy mechanisms such as
government-backed innovation funding (e.g., the UK Net Zero
Innovation Portfolio) could bridge the gap for a demonstration
plant, but durable viability calls for aggregating nitrate feedstock
from multiple treatment works toward larger scales or remunerating
the nitrate-removal service itself, rather than supporting
electricity prices or carbon credits.

The ecological assessment yields a favorable but
accounting-sensitive result. With the fully wind-powered designs
selected by the optimizer, the process is a net consumer of \ce{CO2}
within the gate-to-gate boundary ($-0.73$~kg\ce{CO2} per kg of urea,
the carbon fixed in the product) while additionally removing nitrate
from the water cycle. When the life-cycle intensity of wind
electricity is included as a sensitivity, the high specific
electricity demand of about 190~kWh per kilogram of urea turns the
balance to $+1.6$~kg\ce{CO2}\,kg$^{-1}$, comparable to conventional
production. Reducing the excess-hydrogen operation of the nitrate
hydrogenation and the thermal duties of the ammonia section is
therefore simultaneously the key energy, cost, and emissions lever,
and would make the carbon benefit robust to the accounting choice.

Overall, the proposed process indicates that decentralized,
small-scale urea production from wastewater is technically feasible
and fully compatible with islanded renewable operation, but that its
practical viability is governed by the capital and energy intensity
of small-scale synthesis rather than by renewable intermittency.
Future work should focus on reducing the specific energy demand of
the nitrate-to-ammonia train, exploring process intensification and
modularization to cut the fixed capital of the synthesis sections,
and refining techno-economic assumptions at larger, aggregated
scales. Additionally, more comprehensive life-cycle assessments,
including downstream product utilization and transportation effects,
would further clarify the environmental impact of the system.

\section*{Acknowledgements}

The authors thank the Department of Chemical and Biological
Engineering at Ko\c{c} University for supporting this work.

\section*{Data Availability}

The tools, models, and code used in this work are provided in the
following GitHub repository:
\url{https://github.com/CanberkDemir/Green-Urea}. This repository
contains the computational resources supporting the methodology and
analysis presented in this study.

\section*{Supporting Information}

The following material is available in the Supporting Information:
\begin{itemize}
  \item Screening of water-electrolyzer technologies and derivation
    of the specific energy demand $e_{\mathrm{spec}}$
  \item Tuned kinetic parameters for the nitrate-hydrogenation and
    urea-synthesis reaction networks
  \item Full CAPEX and OPEX formulations used in the NPW analysis
  \item Comparison of the ReLU surrogates with piecewise-linear
    hyperplane-tree surrogates (ridge and elastic-net leaf models)
\end{itemize}

\printbibliography

\newpage

%%%%%%%%%%%%%%%%%%%%%%%%%%%%%%%%%%%%%%%%%%%%%%%%%%%%%%%%%%%%%%%%%%%%%
%% Graphical table-of-contents entry (sized per journal requirements)
%%%%%%%%%%%%%%%%%%%%%%%%%%%%%%%%%%%%%%%%%%%%%%%%%%%%%%%%%%%%%%%%%%%%%
\rule{0.05in}{1.75in}%
\begin{minipage}[b][1.75in]{3.25in}
  \sffamily
  \frenchspacing
  \centering
  % Suggested TOC graphic: a condensed version of the process flow
  % diagram (nitrate brine -> NH3 -> urea/DEF) with the wind/storage
  % loop, e.g. a cropped prodScheme.png.
  %\includegraphics[width=3.1in]{prism-uploads/prodScheme.png}
  Nitrate-rich brine + biogenic \ce{CO2} + offshore wind
  $\rightarrow$ green urea (DEF)
\end{minipage}%
\rule{0.05in}{1.75in}

\end{document}
